%------------------------------------------------------------------------------%
\documentclass[a4paper,12pt,fleqn]{article}

\usepackage{gaal}

\ShowAnswers

%------------------------------------------------------------------------------%
\begin{document}

\makeHeader{GAAL}{Prova 2 -- S}{2026-2}

% 01
%------------------------------------------------------------------------------%
\exercise{20\%}
Determine o volume do paralelogramo determinado pelos vetores
\[
  \qquad
  \vec{p} = \vetortri{1}{2}{-1}
  \qquad
  \vec{q} = \vetortri{2}{-1}{3}
  \qquad
  \vec{r} = \vetortri{0}{4}{1}
\]

\begin{answer}
  \begin{align*}
    p
    & = \left(\vec{p} \times \vec{q}\right) \cdot \vec{r} \\
    & = \begin{vmatrix}
          1 & 2  & -1 \\
          2 & -1 & 3  \\
          0 & 4  & 1
        \end{vmatrix}
    \\
    & = 1 \begin{vmatrix}
        -1 & 3 \\
        4  & 1
      \end{vmatrix}
      - 2 \begin{vmatrix}
        2 & 3 \\
        0 & 1
      \end{vmatrix}
      - 1 \begin{vmatrix}
        2  & -1 \\
        0  & 4
      \end{vmatrix}
    \\
    & = (-1-12) - 2(2) - (8)
    \\
    & = -13 - 4 - 8
    \\
    & = -25
  \end{align*}
  Volume
  \[
    V = \abs{p} = 25
  \]
\end{answer}

% 02
%------------------------------------------------------------------------------%
\exercise{20\%}
Determine o ângulo entre os vetores
\(
  \vec{v} = \vetortri{2}{-1}{2}
\) e
\(
  \vec{w} = \vetortri{-1}{3}{2}
\)

\begin{answer}
  \begin{align*}
    \vec{v} \cdot \vec{w}
    & = \vetortri{2}{-1}{2} \cdot \vetortri{-1}{3}{2}
    \\
    & = 2(-1) + (-1)(3) + 2(2)
    \\
    & = -2 - 3 + 4
    \\
    & = -1
  \end{align*}

  \begin{align*}
    \norm{\vec{v}}
    & = \sqrt{2^2 + (-1)^2 + 2^2}
    \\
    & = \sqrt{4+1+4}
    \\
    & = 3
  \end{align*}

  \begin{align*}
    \norm{\vec{w}}
    & = \sqrt{(-1)^2 + 3^2 + 2^2}
    \\
    & = \sqrt{1+9+4}
    \\
    & = \sqrt{14}
  \end{align*}

  \begin{align*}
    \cos(\theta)
    & = \frac{\vec{v}\cdot\vec{w}}
             {\norm{\vec{v}}\norm{\vec{w}}}
    \\
    & = \frac{-1}{3\sqrt{14}}
  \end{align*}

  \[
    \theta = \arccos\left(-\frac{1}{3\sqrt{14}}\right)
  \]

  \[
    {\color{gray}
      \theta \approx \ang{95.11}
    }
  \]
\end{answer}

% 03
%------------------------------------------------------------------------------%
\exercise{20\%}
Determine a norma do vetor
\(
  \vec{\psi}
  = 2\vetortri{1}{-2}{3}
  - 3\vetortri{2}{1}{-1}
  + 2\vetortri{-1}{0}{2}
\)

\begin{answer}
  \begin{align*}
    \vec{\psi}
    & = 2\vetortri{1}{-2}{3}
      - 3\vetortri{2}{1}{-1}
      + 2\vetortri{-1}{0}{2}
    \\
    & = \vetortri{2}{-4}{6}
      + \vetortri{-6}{-3}{3}
      + \vetortri{-2}{0}{4}
    \\
    & = \vetortri{-6}{-7}{13}
  \end{align*}

  \begin{align*}
    \norm{\vec{\psi}}
    & = \sqrt{(-6)^2+(-7)^2+13^2}
    \\
    & = \sqrt{36+49+169}
    \\
    & = \sqrt{254}
  \end{align*}
\end{answer}

% 04
%------------------------------------------------------------------------------%
\exercise{20\%}
Calcule o determinante da matriz
\[
  A =
  \begin{pmatrix}
    3 & 1 & 0 & 1 \\
    1 & 3 & 1 & 0 \\
    0 & 1 & 3 & 1 \\
    1 & 0 & 1 & 3
  \end{pmatrix}
\]
e determine se a matriz é invertível

\begin{answer}
  Expandindo o determinante pela primeira linha

  \begin{align*}
    \det(A)
    & = 3M_{11} - M_{12} - M_{14}
  \end{align*}

  Calculando os menores

  \begin{align*}
    M_{11}
    & = \begin{vmatrix}
          3 & 1 & 0 \\
          1 & 3 & 1 \\
          0 & 1 & 3
        \end{vmatrix}
    \\
    & = 3(9-1) - 1(3)
    \\
    & = 24-3
    \\
    & = 21
  \end{align*}

  \begin{align*}
    M_{12}
    & = \begin{vmatrix}
          1 & 1 & 0 \\
          0 & 3 & 1 \\
          1 & 1 & 3
        \end{vmatrix}
    \\
    & = 1(9-1) - 1(0-3)
    \\
    & = 8+1
    \\
    & = 9
  \end{align*}

  \begin{align*}
    M_{14}
    & = \begin{vmatrix}
          1 & 3 & 1 \\
          0 & 1 & 3 \\
          1 & 0 & 1
        \end{vmatrix}
    \\
    & = 1(1) - 3(0-3) + 1(0-1)
    \\
    & = 1+9-1
    \\
    & = 9
  \end{align*}

  Portanto

  \begin{align*}
    \det(A)
    & = 3(21)-9-9
    \\
    & = 63-18
    \\
    & = 45
  \end{align*}

  Como
  \[
    \det(A) = 45 \neq 0
  \]
  a matriz $A$ é invertível
\end{answer}

% 05
%------------------------------------------------------------------------------%
\exercise{20\%}
Encontre o vetor unitário na direção e sentido de
\(
  \vec{p}\times\vec{q}
\), sendo
\(
  \vec{p} = \vetortri{1}{2}{2}
\) e
\(
  \vec{q} = \vetortri{2}{-1}{1}
\)

\begin{answer}
  Inicialmente, calculamos o produto vetorial

  \begin{align*}
    \vec{\psi}
    & = \vec{p}\times\vec{q}
    \\
    & = \begin{vmatrix}
          \vec{i} & \vec{j} & \vec{k} \\
          1       & 2       & 2       \\
          2       & -1      & 1
        \end{vmatrix}
    \\
    & = \vec{i}
        \begin{vmatrix}
          2 & 2 \\
          -1 & 1
        \end{vmatrix}
      - \vec{j}
        \begin{vmatrix}
          1 & 2 \\
          2 & 1
        \end{vmatrix}
      + \vec{k}
        \begin{vmatrix}
          1 & 2 \\
          2 & -1
        \end{vmatrix}
    \\
    & = \vec{i}(2+2)
      - \vec{j}(1-4)
      + \vec{k}(-1-4)
    \\
    & = 4\vec{i}+3\vec{j}-5\vec{k}
    \\
    & = \vetortri{4}{3}{-5}
  \end{align*}

  Calculamos agora a norma do vetor

  \begin{align*}
    \norm{\vec{\psi}}
    & = \sqrt{4^2+3^2+(-5)^2}
    \\
    & = \sqrt{16+9+25}
    \\
    & = \sqrt{50}
    \\
    & = 5\sqrt{2}
  \end{align*}

  O vetor unitário procurado é

  \begin{align*}
    \hat{\psi}
    & = \frac{\vec{\psi}}{\norm{\vec{\psi}}}
    \\
    & = \frac{1}{5\sqrt{2}}\vetortri{4}{3}{-5}
    \\
    & = \Vetortri{
          \frac{4}{5\sqrt{2}}
        }{
          \frac{3}{5\sqrt{2}}
        }{
          \frac{-5}{5\sqrt{2}}
        }
  \end{align*}
\end{answer}

%------------------------------------------------------------------------------%
\end{document}
%------------------------------------------------------------------------------%
